\centerline{\bf \Huge Вариант ЕГЭ 2025. Калмыкия}
\centerline{\bf \Huge }

\problem{\textbf{13}} а) Решите уравнение

$2 - 2 \cos^2{x} + \sqrt{2} \sin{x} = \sqrt{2} - 2\sin{(x-\pi)}$

б) Найдите все корни этого уравнения, принадлежащие отрезку $[-2\pi;-\dfrac{\pi}{2}]$.

\problem{\textbf{13*}} а) Решите уравнение $\log _2\left(x^2-2 x\right)=3$.

б) Найдите все корни этого уравнения, принадлежащие отрезку $\left[\log _2 0,1; \log _2 13\right]$.

\problem{\textbf{15}} Решите неравенство $\dfrac{x^3-x^2-x+1}{4^{x^2}-16\cdot2^{x^2}+64}\leqslant 0$

\problem{\textbf{16}} 15 декабря 2026 года планируется взять кредит в банке на сумму 9 млн рублей на 36 месяцев. Условия возврата таковы:

- 1 -го числа каждого месяца долг возрастает на $r$ процентов по сравнению с концом предыдущего месяца;

- co 2-го по 14-е число каждого месяца необходимо одним платежом оплатить часть долга;

- 15-го числа каждого месяца долг должен быть на одну и ту же величину меньше долга на 15-е число предыдущего месяца;

- к 15 декабря 2029 года кредит должен быть полностью погашен.

Чему равно $r$, если общая сумма платежей в 2027 году составит 4830 тыс. рублей?

\problem{\textbf{16*}} Зависимость количества $Q$ в шт. при условии $0 \leqslant Q \leqslant 15000$ купленного у фирмы товара от цены $P$ в руб. за шт. выражается формулой $Q=15000-P$. Затраты на производство $Q$ единиц товара составляют $3000 Q+1000000$ рублей. Кроме затрат на производство, фирма должна платить налог $t$ рублей при условии $0<t<10000$ с каждой произведённой единицы товара. Таким образом, прибыль фирмы составляет $P Q-3000 Q-1000000-t Q$ рублей, а общая сумма налогов, собранных государством, равна $t Q$ рублей.
Фирма производит такое количество товара, при котором её прибыль максимальна. При каком значении $t$ общая сумма налогов, собранных государством, будет максимальной?

\newpage
\hspace{3mm}

\problem{\textbf{17}} В четырёхугольник $KLMN$ вписана окружность с центром $O$. Эта окружность касается стороны $MN$ в точке $A$. Известно, что $\angle MNK=90°$,  $\angle NKL=\angle KLM=120°$.

а) Докажите, что точка $A$ лежит на прямой $LO$.

б) Найдите длину стороны $MN$, если $LA=1$.

\problem{\textbf{17*}} Дан ромб $A B C D$. На диагонали $A C$ отмечены точки $M$ и $N$, так что $A M=N M=N C$. Прямая $B M$ пересекает сторону $A D$ в точке $P$, а прямая $B N$ пересекает сторону $C D$ в точке $Q$.

а) Докажите, что площадь четырехугольника $B P D Q$ равна площади треугольника $A D C$.

б) Найдите $B D$, если известно, что $A C=2 \sqrt{3}$ и около пятиугольника $M N Q D P$ можно описать окружность.

\problem{\textbf{19*}} На доске написано 30 чисел: десять «5», десять «4» и десять «3». Эти числа разбивают на две группы, в каждой из которых есть хотя бы одно число. Среднее арифметическое чисел в первой группе равно $A$, среднее арифметическое чисел во второй группе равно $B$. При этом для группы из единственного числа среднее арифметическое равно этому числу.

a) Приведите пример разбиения исходных чисел на две группы, при котором среднее арифметическое всех чисел меньше $\dfrac{A+B}{2}$.

б) Докажите, что если разбить исходные числа на две группы по 15 чисел, то среднее арифметическое всех чисел будет равно $\dfrac{A+B}{2}$.

в) Найдите наибольшее возможное значение выражения $\dfrac{A+B}{2}$.